\documentclass[12pt, a4paper]{article}

\usepackage[utf8]{inputenc}
\usepackage{geometry}
\usepackage{listings}
\usepackage{xcolor}
\usepackage{hyperref}
\usepackage{graphicx}
\usepackage{booktabs}
\usepackage{amsmath}

\geometry{top=1in, bottom=1in, left=1in, right=1in}

\definecolor{codegreen}{rgb}{0,0.6,0}
\definecolor{codegray}{rgb}{0.5,0.5,0.5}
\definecolor{codepurple}{rgb}{0.58,0,0.82}
\definecolor{backcolour}{rgb}{0.95,0.95,0.92}

\lstdefinestyle{mystyle}{
    backgroundcolor=\color{backcolour},   
    commentstyle=\color{codegreen},
    keywordstyle=\color{magenta},
    numberstyle=\tiny\color{codegray},
    stringstyle=\color{codepurple},
    basicstyle=\ttfamily\footnotesize,
    breakatwhitespace=false,         
    breaklines=true,                 
    captionpos=b,                    
    keepspaces=true,                 
    numbers=left,                    
    numbersep=5pt,                  
    showspaces=false,                
    showstringspaces=false,
    showtabs=false,                  
    tabsize=2
}
\lstset{style=mystyle}

\title{\textbf{Java-Safe C++ (jscpp)}\\ \large A C++-Subset Compiler with a Bytecode VM for Portable, Memory-Safe Execution}
\author{Engineering Team}
\date{\today}

\begin{document}

\maketitle

\begin{abstract}
C++ is a powerful systems programming language but notoriously lacks memory safety, often leading to undefined behavior, buffer overflows, and null pointer dereferences. This project, \textbf{Java-Safe C++ (jscpp)}, is an end-to-end compiler and stack-based virtual machine built entirely from scratch. It compiles a restricted subset of C++ into a custom portable bytecode format and executes it within a managed Virtual Machine. The compiler automatically injects safety checks (such as null pointer and array bounds checking) and manages memory through Automatic Reference Counting (ARC), effectively bringing Java-like safety guarantees to C++ syntax.
\end{abstract}

\tableofcontents
\newpage

\section{Introduction}
The primary objective of this project is to demonstrate modern compiler construction principles combined with managed runtime execution. Instead of relying on existing frameworks like LLVM or interpreters like Python, we implemented a full multi-stage compiler pipeline in standard C++17. 

\subsection{Design Goals}
\begin{itemize}
    \item \textbf{Memory Safety:} Unsupported or dangerous operations result in graceful runtime exceptions rather than host OS segmentation faults.
    \item \textbf{No Undefined Behavior:} Static typing and runtime safety eliminate undefined behavior for the supported subset.
    \item \textbf{Zero External Dependencies:} The project uses only the standard C++ library (no Lex/Yacc, no LLVM).
    \item \textbf{Understandable Architecture:} Clean separation of Lexical Analysis, Parsing, Semantic Analysis, Code Generation, and VM execution.
\end{itemize}

\section{Language Specification}
The compiler supports a minimal but Turing-complete subset of C++:

\subsection{Supported Features}
\begin{itemize}
    \item \textbf{Types:} \texttt{int}, \texttt{bool}, \texttt{void}, Pointers (\texttt{T*}), and 1D Arrays (\texttt{T[N]}).
    \item \textbf{Control Flow:} \texttt{if / else}, \texttt{while}, \texttt{for}, \texttt{return}.
    \item \textbf{Expressions:} Arithmetic (\texttt{+}, \texttt{-}, \texttt{*}, \texttt{/}, \texttt{\%}) with 32-bit wrap-around, relational (\texttt{==}, \texttt{<}, \texttt{>=}, etc.), short-circuit logical operators (\texttt{\&\&}, \texttt{||}, \texttt{!}), compound assignment (\texttt{+=}, \texttt{-=}, \ldots) and \texttt{++}/\texttt{--}.
    \item \textbf{Memory:} Dynamic allocation via \texttt{new int} / \texttt{new bool}, pointer dereferencing (\texttt{*p}), array indexing (\texttt{arr[i]}), and \texttt{delete}.
    \item \textbf{Functions:} User-defined functions with parameters, defined in any order and possibly recursive, plus a built-in \texttt{print()} function.
\end{itemize}

\subsection{Unsupported Features}
To keep the project focused on safety mechanisms, we explicitly reject features like templates, classes/structs, standard library (STL) headers, preprocessor macros, and multiple inheritance. Constructs that could create an unsafe pointer are also rejected at compile time, each with its own error message: the address-of operator \texttt{\&}, pointer arithmetic, pointers to pointers, arrays of pointers, global variables, whole-array assignment, and non-void functions that can end without returning a value.

\section{Compiler Pipeline Architecture}
The compiler is divided into five distinct stages:

\subsection{1. Lexer}
Converts the raw source code string into a sequence of \texttt{Token} objects. It handles whitespace/comment stripping and precise line/column source-location tracking for detailed error diagnostics.

\subsection{2. Parser}
A handcrafted recursive-descent parser that consumes tokens and builds a custom Abstract Syntax Tree (AST). It properly enforces operator precedence (e.g., \texttt{a + b * c} parses as \texttt{a + (b * c)}).

\subsection{3. Semantic Analyzer}
Implemented as an AST Visitor, this phase maintains lexical scopes and symbol tables. It runs in two passes: the first records the signature of every function (so functions may call each other in any order), and the second type-checks every body. It checks that the left side of an assignment is a real storage location and that every path of a non-void function returns. Any violation results in a compile-time \texttt{SemanticError}, printed with the offending source line and a caret.

\subsection{4. Code Generator}
Lowers the AST into our custom bytecode. Crucially, this is where memory safety is enforced. The generator automatically injects instructions like \texttt{CHECK\_NULL}, \texttt{CHECK\_BOUNDS}, \texttt{INC\_REF}, and \texttt{DEC\_REF} around memory access and pointer assignments.

\subsection{5. Virtual Machine}
A deterministically evaluated, stack-based runtime engine. The VM maintains an operand stack, a call stack for function frames, local variables, and a managed heap. Every stack, local-variable and heap access is range-checked, and a runtime exception reports the source line and function and makes \texttt{jscpp} exit with status 1.

\section{Memory Safety Mechanisms}
The core achievement of \textbf{jscpp} is genuine safety injection. The compiler protects the program using three primary runtime guards:

\subsection{Null Pointer Checks}
Before any pointer dereference (\texttt{*p}), the compiler emits a \texttt{CHECK\_NULL} opcode. If the pointer evaluates to address \texttt{0}, the VM halts and throws a \texttt{NullPointerException}.

\subsection{Array Bounds Checks}
When accessing arrays (\texttt{arr[i]}), the compiler pushes the array pointer and the index, then emits \texttt{CHECK\_BOUNDS}. The VM's managed heap stores the allocated size of every object. If $i < 0$ or $i \ge \text{size}$, an \texttt{ArrayIndexOutOfBoundsException} is thrown.

\subsection{Automatic Reference Counting (ARC)}
All \texttt{new} allocations and arrays live on the VM's managed heap with a reference count. Every local pointer or array variable owns one reference. Values produced by \texttt{new} or by a function returning a pointer are \emph{owned} (already counted), so storing them transfers the reference; values read from a variable are \emph{borrowed}, so storing them emits \texttt{INC\_REF}. The compiler emits \texttt{DEC\_REF} when a variable is reassigned, at the end of every block, and on every \texttt{return}. Pointer arguments become owned by the callee, which releases them on return. When a count reaches zero the object is freed immediately. Running with \texttt{--stats} prints the number of allocations and frees; for every example that finishes normally, no objects are alive at exit.

\subsection{Safe \texttt{delete}}
\texttt{delete p} releases \texttt{p}'s reference and sets \texttt{p} to \texttt{nullptr}. Another variable that refers to the same object keeps it alive, so \texttt{delete} can never create a dangling pointer or a double free.

\subsection{Defined Integer Arithmetic}
Signed overflow is undefined behaviour in C++. The VM therefore computes \texttt{+}, \texttt{-} and \texttt{*} on unsigned 32-bit values and converts back, so results wrap exactly like Java's \texttt{int}. \texttt{INT\_MIN / -1} and division by zero are handled explicitly.

\section{Bytecode Format}
The compiled bytecode can be executed in memory or serialized to a binary \texttt{.jbc} file. The file stores every integer in little-endian order, \texttt{CALL} refers to functions by their index in a function table, and each instruction carries its source line. A loaded file is verified (opcodes, constant and function indices, jump targets, allocation sizes) before it runs, and a file compiled on Windows runs unchanged on Linux. The instruction set includes operations such as:
\begin{itemize}
    \item \texttt{PUSH\_CONST}, \texttt{LOAD\_LOCAL}, \texttt{STORE\_LOCAL}
    \item Arithmetic \& Logical operators (\texttt{ADD}, \texttt{MUL}, \texttt{AND}, \texttt{EQ}, etc.)
    \item Control Flow (\texttt{JUMP}, \texttt{JUMP\_IF\_FALSE}, \texttt{CALL}, \texttt{RETURN})
    \item Memory \& Safety (\texttt{ALLOC}, \texttt{CHECK\_BOUNDS}, \texttt{CHECK\_NULL}, \texttt{INC\_REF}, \texttt{DEC\_REF}, \texttt{LOAD\_ARRAY}, \texttt{STORE\_ARRAY\_INV})
\end{itemize}

\section{Build and Execution Instructions}

\subsection{Building the Compiler}
The project uses CMake and requires a C++17 compatible compiler.
\begin{lstlisting}[language=bash]
cd java-safe-cpp
mkdir build
cd build
cmake ..
cmake --build .
\end{lstlisting}

\subsection{Command Line Interface (CLI)}
The compiled \texttt{jscpp} executable provides several commands:
\begin{lstlisting}[language=bash]
# Run a program directly from source
./jscpp run examples/arrays.cppsafe

# Run and print heap allocation / free counts
./jscpp run examples/ownership.cppsafe --stats

# View VM execution trace with the operand stack
./jscpp run examples/pointers.cppsafe --trace

# Show the output of the lexer and the parser
./jscpp tokens examples/hello.cppsafe
./jscpp ast examples/refcount.cppsafe

# Compile to binary bytecode, run it, and dump it
./jscpp compile examples/hello.cppsafe -o hello.jbc
./jscpp run hello.jbc
./jscpp dump hello.jbc
\end{lstlisting}

\section{Examples and Expected Output}

\subsection{Array Bounds Safety}
\textbf{Source Code:}
\begin{lstlisting}[language=C++]
void main() {
    int arr[3];
    arr[0] = 10;
    print(arr[0]);
    arr[5] = 100; // Trigger bounds check
}
\end{lstlisting}
\textbf{Output:}
\begin{lstlisting}
10
Runtime Error:
ArrayIndexOutOfBoundsException: index 5, length 3
    at line 7 in main() (instruction 20)
\end{lstlisting}

\subsection{Null Pointer Safety}
\textbf{Source Code:}
\begin{lstlisting}[language=C++]
void main() {
    int* p = nullptr;
    print(*p); // Trigger null check
}
\end{lstlisting}
\textbf{Output:}
\begin{lstlisting}
Runtime Error:
NullPointerException: attempted to dereference a null pointer
    at line 5 in main() (instruction 6)
\end{lstlisting}

\section{Faults Found and Fixed}
The first working version was audited by building it and running small probe programs. The most serious faults were:
\begin{itemize}
    \item \texttt{\&\&} and \texttt{||} emitted a jump to address 0 that was never patched, so any condition using them looped forever.
    \item Parameters were numbered after the locals of earlier functions, so a call stored its arguments in the wrong slots (a program that should print 42 printed 0).
    \item Variables declared at function level were never released, so reference counting freed nothing; arrays were never released either.
    \item Functions could not be called before their definition, any expression was accepted on the left of \texttt{=}, and global variables compiled but read the wrong memory.
    \item Signed overflow in the VM was undefined behaviour in the host C++, runtime errors exited with status 0, \texttt{CALL} relied on \texttt{std::hash} (not portable), and loaded bytecode was not checked.
\end{itemize}
The full list of twenty faults, how each was confirmed, and the fix plan are in \texttt{docs/audit-and-plan.md}.

\section{Testing}
The project has an automated test suite driven by CTest (\texttt{ctest --test-dir build}). Every program in \texttt{examples/} is run and its exact output, including the heap statistics line, is compared with a stored expected result; for the programs that demonstrate a runtime exception, the exception text and source line are also checked. Each example is run twice: from source, and after being saved to and loaded from a \texttt{.jbc} file. Another 28 small programs must each be rejected with a specific compile error. All 62 tests pass on Windows (MinGW g++ 15.2) and on Linux (Ubuntu under WSL, g++ 15.2), and the build produces no warnings with \texttt{-Wall -Wextra -Wpedantic}.

\section{Conclusion and Limitations}
This project demonstrates how memory safety features can be built into a compiler for C++-style syntax. By tightly coupling a custom compiler pipeline with a strict managed runtime, we achieve a memory-safe environment that prevents undefined behavior for our specific language subset.

\textbf{Disclaimer on Portability and Scope:} It is important to emphasize that this project does \emph{not} claim to make C++ "Write Once, Run Anywhere" (WORA). The portability and safety guarantees apply strictly to the minimal academic subset of C++ implemented by our custom VM. Real-world C++ constructs such as templates, classes/structs, STL, preprocessor macros, and multiple inheritance are explicitly excluded from this compiler.

\end{document}
